
On
trouve donc un homomorphisme de spécialisation sur \(K_\bullet(X_t)\), mais à valeurs non dans \(K_\bullet(X_s)\), mais dans \(K_\bullet(X_{s'})\), où \(X_{s'}\) est déduit de \(X_s\) par une extension \(k(s)\to k(s')\) des corps de base (qu'on peut supposer extension de type fini si \(S\) est universellement japonais). Passant aux fibres géométriques comme expliqué dans 7.16, on trouve de même un homomorphisme
\begin{equation}\tag{7.17.2}
  \sigma_\bullet:K_\bullet(X_{\bar t})\longrightarrow K_\bullet(X_{\overline{s'}});
\end{equation}
donc ici encore, l'homomorphisme de spécialisation «~géométrique~» prend ses valeurs non dans \(K_\bullet(X_{\bar s})\), mais dans \(K_\bullet(X_{\overline{s'}})\), où \(X_{\overline{s'}}\) est déduit de \(X_{\bar s}\) par l'extension de corps algébriquement clos \(k(\bar s)\to k(\overline{s'})\).

\textbf{7.17.3.} En général, l'homomorphisme \(K_\bullet(X_{\bar s})\to K_\bullet(X_{\overline{s'}})\) n'est pas un isomorphisme (à cause de la présence possible de composantes «~continues~» dans le groupe \(K_\bullet(X_{\bar s})\)). Cependant, il est facile de vérifier que c'est un isomorphisme modulo équivalence algébrique (7.11), et aussi modulo équivalence numérique (4.5.1) si \(f\) est lisse et propre. Par suite, utilisant 7.12, on trouve que (7.17.2) induit un homomorphisme de spécialisation
\begin{equation}\tag{7.17.3.1}
  K_\bullet^{\mathrm{alg}}(X_{\bar t})\longrightarrow
  K_\bullet^{\mathrm{alg}}(X_{\bar s});
\end{equation}
et si \(f\) est propre et lisse, et sous réserve que l'homomorphisme de spécialisation (7.17.2) soit compatible avec l'équivalence numérique (cf. 7.9.4), on trouve de même un homomorphisme
\begin{equation}\tag{7.17.3.2}
  K_\bullet^{\mathrm{num}}(X_{\bar t})\longrightarrow
  K_\bullet^{\mathrm{num}}(X_{\bar s}).
\end{equation}
Nous laissons au lecteur le soin de développer (en cas de besoin) les constructions correspondantes pour \(\Gr_\bullet\), \(K^\bullet\), et \(\Gr^\bullet\), et les propriétés de ces homomorphismes de spécialisation qui résultent trivialement des propriétés développées dans les sections précédentes. Signalons seulement
